(1) 当地竖直向上的地球自转分量 $\Omega_z=\Omega_E\sin\varphi$。旋转坐标中科里奥利加速度为 $-2\boldsymbol\Omega\times\boldsymbol v$，水平分量给 $\ddot x+\omega_0^2x=2\Omega_z\dot y$、$\ddot y+\omega_0^2y=-2\Omega_z\dot x$。横向自转分量对水平速度的科里奥利加速度沿竖直，已在理想水平近似下去掉。
(2) 定义 $z=x+iy$，两式相加得到 $\ddot z+2i\Omega_z\dot z+\omega_0^2z=0$。代入 $z\sim e^{-i\omega t}$ 后有 $\omega=\Omega_z\pm\sqrt{\omega_0^2+\Omega_z^2}$。对 $|\Omega_z|\ll\omega_0$，快振动与慢转向叠加，即线摆近似 $z(t)=e^{-i\Omega_zt}[A\cos(\omega_0t)+B\sin(\omega_0t)]$，其中 $A,B$ 为适当复常数。
(3) 线摆主轴相对当地地面每秒转角约为 $-\Omega_z=-\Omega_E\sin\varphi$，北半球从上方观察为顺时针，南半球为逆时针，赤道一阶无进动。对于非赤道纬度，摆面方向转过 $2\pi$ 的时间为 $P_F=2\pi/(\Omega_E|\sin\varphi|)$；进动只依赖当地地球自转竖直分量而非摆长。
(4) 若两个水平分量等幅且相差四分之一周期，轨迹可为旋转的圆而没有唯一线摆主轴；不能将圆周的快角速度误判为福科摆面慢进动。方程和近似仅保持一阶 $\Omega_E/\omega_0$ 的水平科里奥利效应，忽略地球曲率、椭圆偏振各向异性、悬点漂移、空气摩擦和大角度非线性。
